%! Vector-Valued Functions — Unit 4, Lesson 4.9 — Homework
\documentclass[10pt]{article}
\usepackage{precalculus-article}
\usepackage{precalculus-boxes}
\newcommand{\TallMath}[1]{$\displaystyle #1\rule[-1.4em]{0pt}{3.2em}$}

\begin{document}

\pageheader{Unit 4, Lesson 4.9}{Homework}
\namedateperiod

\vspace{0.12in}

\begin{notesbox}{Vector-Valued Functions Practice}
\small

\textbf{1.}\quad A particle's position is given by the parametric function
$f(t)=(5t-2,\;t^{2}+1)$.

\begin{enumerate}[label=(\alph*)]
  \item Rewrite $f(t)$ as a vector-valued position function $\vec{p}(t)$.

  \vspace{0.06in}
  \hspace{1em} $\vec{p}(t) = \langle \blank{5cm}\rangle$

  \vspace{0.1in}
  \item Find $\vec{p}(3)$.

  \vspace{0.06in}
  \hspace{1em} $\vec{p}(3) = \langle \blank{4cm}\rangle$
\end{enumerate}

\vspace{0.14in}

\textbf{2.}\quad Using your answer from Problem 1, find $\|\vec{p}(3)\|$.
Round to the nearest tenth.

\vspace{0.06in}
\hspace{1em} $\|\vec{p}(3)\| = \sqrt{\blank{1.8cm}^{2}+\blank{1.8cm}^{2}}
= \sqrt{\blank{2.5cm}} \approx $ \blank{2.5cm}

\vspace{0.14in}

\textbf{3.}\quad A particle's velocity vector at a certain instant is
$\vec{v}=\langle -3,\;5\rangle$.

\begin{enumerate}[label=(\alph*)]
  \item Is the particle moving right or left?  Explain using the sign of the
        $x$-component.

  \writeline

  \item Is the particle moving up or down?  Explain using the sign of the
        $y$-component.

  \writeline
\end{enumerate}

\vspace{0.14in}

\textbf{4.}\quad At time $t=2$ a particle has $x'(2)=6$ and $y'(2)=-8$.
Find the speed at $t=2$.

\vspace{0.06in}
\hspace{1em} $\vec{v}(2) = \langle \blank{3cm}\rangle$ \qquad
Speed $= \sqrt{\blank{2.5cm}} = $ \blank{3cm}

\vspace{0.14in}

\textbf{5.}\quad \textbf{(AP-Style Multiple Choice)}\quad
A particle moves with position vector
$\vec{p}(t)=\langle 4\cos t,\;3\sin t\rangle$ for $0\le t\le 2\pi$.

At which value of $t$ is the particle \emph{farthest} from the origin?

\begin{enumerate}[label=(\Alph*)]
  \item $t=0$
  \item $t=\dfrac{\pi}{2}$
  \item $t=\pi$
  \item The distance from the origin is the same for all $t$.
\end{enumerate}

\end{notesbox}

\end{document}
